\documentclass[12pt,a4paper,oneside]{article}
\usepackage[brazil]{babel}
\usepackage[utf8]{inputenc}
\usepackage{olymp}

\setcounter{problem}{0}

\begin{document}

\begin{problem}{Extrato bancário}{extrato}{2}

Preocupado com os gastos do mês, Manuel retirou um extrato da sua conta bancária e agora quer saber quanto entrou e quanto saiu da sua conta.
O extrato mostra os valores de todas as movimentações, incluindo depósitos, transferências, pix, pagamentos de contas, entre outras.
Todos os valores positivos são valores que entraram na sua conta e todos os negativos são valores que saíram da conta. Faça um programa para totalizar esses valores.

%\Notes
%
%Observações, se tiver.

\InputFile

A entrada começa com uma linha contendo um inteiro $N$, o número de movimentações na conta.
Em seguida há $N$ linhas, cada uma com o valor real $V$ de uma movimentação realizada.
Restrições: $1 \leq N \leq 100$ e $-10000 \leq V \leq 10000$.

\OutputFile

Seu programa deve gerar duas linhas na saída: a primeira com o valor total que entrou na conta e a segunda com o valor total que saiu da conta.
Ambos devem ser formatados para duas casas decimais.

%\Notes

\Example

\begin{example}
\exmp{
5
100.00
100.00
-50.00
100.00
-50.00
}{
Entrou: 300.00
Saiu: 100.00
}%
\end{example}

\begin{example}
\exmp{
3
-100.00
-200.00
-300.00
}{
Entrou: 0.00
Saiu: 600.00
}%
\end{example}

\begin{example}
\exmp{
4
1.50
2.50
3.50
-5.00
}{
Entrou: 7.50
Saiu: 5.00
}%
\end{example}

%Comentários sobre o exemplo, se necessário.

\end{problem}

\end{document}
